
\chapter{The moonshine module}\label{chap:moonshine}

The Frenkel--Lepowsky--Meurman moonshine module $V^\natural$
\cite{FLM88} is a holomorphic vertex algebra at $c = 24$ with
$\dim V_1^\natural = 0$. It occupies the unique $V_1=0$ slot in the
Schellekens $c=24$ census.
The twenty-four Niemeier lattice VOAs share $c = 24$ and
$\kappa(V_\Lambda) = \operatorname{rank}(\Lambda) = 24$; the
moonshine module has $\kappa(V^\natural) = c/2 = 12$
\textup{(}Virasoro-sector value;
Proposition~\ref{prop:moonshine-kappa}\textup{)}, and the modular
characteristic separates it from all twenty-four lattice VOAs at
the first shadow level.

On the Virasoro stress-tensor projection, the Leech-to-Monster
comparison changes from the Gaussian lattice lane to the class
$\mathbf{M}$ Virasoro lane. This is an instance of the
quadrichotomy $\{\mathbf{G},\mathbf{L},\mathbf{C},\mathbf{M}\}$
analysed in Chapter~\ref{chap:shadow-quadrichotomy-platonic}, not a
new moonshine archetype. The Virasoro-sector complementarity sum
$\kappa+\kappa' = 13$ witnesses the universal Koszul conductor
$K = 13$ for that lane, computed from
$K=-c_{\mathrm{ghost}}(\BRST(\mathrm{Vir}))$ in
Chapter~\ref{chap:universal-conductor}.

\begin{table}[ht]
\centering
\small
\caption{Five theorem packages on the moonshine
module~$V^\natural$.}\label{tab:moonshine-five-theorems}
\begin{tabular}{p{2.4cm}p{8.6cm}p{3.2cm}}
\toprule
\textbf{Package} & \textbf{Moonshine assertion} & \textbf{Anchor} \\
\midrule
A (adjunction) &
 $B(V^\natural)$ is the bar coalgebra in the chiral adjunction. The
 Virasoro scalar lane has Verdier partner $\mathrm{Vir}_2$; the full
 algebra $(V^\natural)^!$ has not been computed beyond the
 Virasoro projection.
 & Rem.~\ref{rem:moonshine-koszul-dual} \\
B (inversion) &
 $\Omega(\barB(V^\natural))\xrightarrow{\sim}V^\natural$ is
 bar--cobar inversion on the completed Koszul/PBW lane, not an
 identification of the Koszul dual algebra.
 & Thm.~\ref{thm:bar-cobar-inversion-qi} \\
C (complementarity) &
 $\kappa(V^\natural)+\kappa(\mathrm{Vir}_2)=12+1=13$ on the
 Virasoro scalar lane.
 & Eq.~\eqref{eq:moonshine-complementarity} \\
D (modular char.) &
 $\kappa(V^\natural)=12$, with no rank-$24$ lattice current
 contribution.
 & Prop.~\ref{prop:moonshine-kappa} \\
H (chiral Hochschild) &
On a smooth projective curve~$X$ satisfying the finite-completion
hypotheses, $Z_{\mathrm{ch},X}^{\mathrm{der}}(V^\natural)
=C^\bullet_{\mathrm{ch},X}(V^\natural,V^\natural)$ is the
cochain-level closed-sector object governed by the $[0,2]$
concentration theorem.
 & Thm.~\ref{thm:hochschild-polynomial-growth} \\
\bottomrule
\end{tabular}
\end{table}

No topological, categorical, product-CY, or Hall Hochschild theory is
used in this table. Each such theory requires its own ambient category
and comparison functor before it can be compared with the curve-level
chiral Hochschild complex.

\begin{table}[ht]
\centering
\small
\caption{Shadow archetype data for the Virasoro scalar projection of
the moonshine module.}\label{tab:moonshine-shadow-archetype}
\begin{tabular}{ll}
\toprule
\textbf{Invariant} & \textbf{Value} \\
\midrule
Class & $\mathbf{M}$ on the Virasoro scalar line \\
Shadow depth $r_{\max}$ & $\infty$ on that line \\
$\kappa(V^\natural)$ & $12$ \\
$S_4^{\mathrm{Vir}}$ & $5/1704$ \\
Critical discriminant $\Delta_T$ & $20/71 \neq 0$ \\
$F_1$ & $1/2$ \\
$F_2$ & $7/480$ \\
$r$-matrix & $r^{\mathrm{Vir}}_{24}(z)=12/z^3+2T/z$ \\
Verdier partner (Virasoro lane) & $\mathrm{Vir}_2$ \\
Complementarity sum & $\kappa+\kappa'=13$ \\
$L_\infty$-formality & $T$-line non-formal, since $\Delta_T \neq 0$ \\
\bottomrule
\end{tabular}
\end{table}

%% ================================================================
%% Section 1. The modular characteristic
%% ================================================================
\section{The modular characteristic}

\begin{proposition}[$\kappa(V^\natural) = 12$;
\ClaimStatusProvedHere]\label{prop:moonshine-kappa}%
\index{Monster module!modular characteristic}%
The modular characteristic of the moonshine module is
\[
\kappa(V^\natural) \;=\; \frac{c}{2} \;=\; 12.
\]
\end{proposition}

\begin{proof}
\emph{Bar complex first principles.}
The genus-$1$ fiberwise curvature
$d_{\mathrm{fib}}^2 = \kappa \cdot \omega_1$ is sourced
additively by each independent curvature channel.
Weight-$1$ currents contribute $\kappa = 1$ per boson
(Theorem~\ref{thm:frame-genus1-curvature}); the Virasoro
sector contributes $\kappa = c/2$
(Theorem~\ref{thm:genus1-universal-curvature}).
Since $\dim V_1^\natural = 0$, no Heisenberg channel exists.
The weight-$2$ Griess algebra generators are conformal
primaries: they enter the bar differential at degree~$\geq 2$
(the shadow coefficients $S_3, S_4, \ldots$) but do not
source the scalar genus-$1$ curvature (the degree-$0$ vacuum
diagram has no external insertions). Therefore
$\kappa(V^\natural) = c/2 = 12$.

\emph{Orbifold from the Leech lattice.}
The FLM construction realizes $V^\natural = V_{\Lambda_{24}}^{\bZ/2}$,
where $\bZ/2$ acts by the Leech involution $v \mapsto -v$ combined
with the canonical automorphism. The Leech lattice VOA has
$\kappa(V_{\Lambda_{24}}) = \operatorname{rank}(\Lambda_{24}) = 24$
(Theorem~\ref{thm:lattice:curvature-braiding-orthogonal}).
That value belongs to the rank-$24$ Heisenberg current sector. The
orbifold fixed-point algebra has no weight-$1$ currents, so the
lattice rank formula is not transported across the orbifold. The
remaining scalar curvature source is the primitive stress tensor,
hence $\kappa(V^\natural)=24/2=12$.

\emph{Normalization check.}
The landscape census records the same values
$\kappa(V_{\Lambda_{24}})=24$ and $\kappa(V^\natural)=12$, and
\texttt{compute/lib/moonshine\_shadow\_depth.py} implements the
same separation through \texttt{niemeier\_kappa} and
\texttt{monster\_kappa}. Thus the alternatives $\kappa=c=24$ and
$\kappa=\operatorname{rank}(\Lambda_{24})=24$ are lattice-lane
values, not moonshine scalar values.
\end{proof}

\begin{remark}[Why $\kappa = c/2$, not $\kappa = 24$]%
\label{rem:moonshine-kappa-not-24}%
\index{Monster module!why kappa is not 24}%
The lattice formula $\kappa = \operatorname{rank}(\Lambda)$ applies
to lattice VOAs because each weight-$1$ Heisenberg boson
contributes $\kappa = 1$ to the genus-$1$ obstruction. Lattice
VOAs at $c = 24$ have $\operatorname{rank} = 24 = c$; this
numerical coincidence is an accident of rank-$24$ lattices, not a
general identity. The Monster module is not a lattice VOA:
$\dim V_1^\natural = 0$ eliminates the Heisenberg sector, and the
stress tensor $T$ is a genuine weight-$2$ generator (not the
Sugawara composite $\tfrac{1}{2}\sum j_a j_a$ of weight-$1$
currents). The formula $\kappa = c/2$ applies because the
Virasoro sector alone sources the genus-$1$ bar curvature.
\end{remark}

%% ================================================================
%% Section 2. Shadow depth and the single-line dichotomy
%% ================================================================
\section{Shadow depth and the single-line dichotomy}

On the $T$-line (the single primary line of conformal weight~$2$),
the shadow metric is
\begin{equation}\label{eq:moonshine-shadow-metric}
Q_T(t)
\;=\;
(2\kappa + 3\alpha t)^2 + 2\Delta\,t^2
\;=\;
(24 + 6t)^2 + \tfrac{40}{71}\,t^2,
\end{equation}
with $\alpha = 2$ (universal for Virasoro),
$S_4 = 10/\bigl(c(5c{+}22)\bigr) = 5/1704$, and
critical discriminant
\begin{equation}\label{eq:moonshine-discriminant}
\Delta
\;=\;
8\kappa S_4
\;=\;
8 \cdot 12 \cdot \frac{5}{1704}
\;=\;
\frac{20}{71}
\;\neq\; 0.
\end{equation}
The single-line dichotomy
(Theorem~\ref{thm:single-line-dichotomy}) forces infinite
shadow depth on $L_T$: $r_{\max}(L_T) = \infty$. Thus the
stress-tensor line of $V^\natural$ is class~M alongside the
Virasoro and specified $\mathcal{W}$ stress-tensor lines. The
shadow obstruction tower on this line does not terminate; the
quartic and all higher shadows are generically nonzero.

\begin{remark}[Niemeier discrimination on the Virasoro scalar projection]%
\label{rem:moonshine-niemeier-discrimination}%
\index{Monster module!Niemeier discrimination}%
\index{Niemeier lattice!shadow discrimination from Monster}%
The triple $(\kappa, \text{class}, \Delta)$ separates
$V^\natural$ from all twenty-four Niemeier lattice VOAs at the
Virasoro scalar-shadow level.
The Niemeier lattice VOAs share
$(\kappa, \text{class}, \Delta) = (24, \mathbf{G}, 0)$:
the shadow obstruction tower terminates at degree~$2$
(Theorem~\ref{thm:lattice:niemeier-shadow-universality}).
The Leech lattice VOA $V_{\Lambda_{24}}$ is
indistinguishable from the other twenty-three Niemeier lattice
VOAs at the shadow level; it is distinguished by its theta
series (zero roots), not by the shadow tower.
The stress-tensor line of the moonshine module has
$(\kappa, \text{class}, \Delta) = (12, \mathbf{M}, 20/71)$:
a different shadow profile, separated already by~$\kappa$.

The genus-$1$ free energies reflect the separation:
\begin{equation}\label{eq:moonshine-f1-comparison}
F_1(V^\natural) = \frac{\kappa}{24} = \frac{1}{2},
\qquad
F_1(V_{\Lambda_{24}}) = \frac{24}{24} = 1.
\end{equation}
At genus~$2$:
$F_2(V^\natural) = 7\kappa/5760 = 7/480$, while
$F_2(V_{\Lambda_{24}}) = 7 \cdot 24/5760 = 7/240$.
The ratio $F_g(V^\natural)/F_g(V_\Lambda) = 1/2$ persists
at all genera on the Virasoro scalar projection, where Theorem~D gives
the coefficient $\kappa\lambda_g$. It is not a statement about the
full Monster-equivariant bar complex: the weight-$2$ Griess algebra
contributes additional non-scalar sectors.
\end{remark}

%% ================================================================
%% Section 3. The Koszul dual and complementarity
%% ================================================================
\section{The Koszul dual and complementarity}

\begin{remark}[The Koszul dual of $V^\natural$]%
\label{rem:moonshine-koszul-dual}%
\index{Monster module!Koszul dual}%
The Virasoro Koszul duality $\mathrm{Vir}_c^! \simeq
\mathrm{Vir}_{26-c}$ (Proposition~\ref{prop:virasoro-generic-koszul-dual})
specializes at $c = 24$ to a Virasoro-sector partner
$\mathrm{Vir}_2$. For the full moonshine module the chain object is
$\barB(V^\natural)$, the intrinsic dual coalgebra is
$(V^\natural)^i=H^*(\barB(V^\natural))$, and a strict algebra
$(V^\natural)^!$ would require Verdier/Koszul duality on that
coalgebra. Beyond the Virasoro sector this bar cohomology has not
been computed. The first obstruction is already the weight-$2$
Griess algebra, whose $196883$ primary directions under the Monster
group $\mathbb{M}$ are not a standard-family quadratic dual.
Consequently this chapter records only the Virasoro-sector partner.
The derived centre is a different curve-level cochain object:
$Z_{\mathrm{ch},X}^{\mathrm{der}}(V^\natural)
=C^\bullet_{\mathrm{ch},X}(V^\natural,V^\natural)$, not
$B(V^\natural)$, $(V^\natural)^i$, or $(V^\natural)^!$; it is not a
physical bulk sector without an additional field-theoretic comparison
(cf.\ Remark~\ref{rem:lattice:monster-shadow}).
\end{remark}

The Virasoro-sector complementarity gives
\begin{equation}\label{eq:moonshine-complementarity}
\kappa(V^\natural) + \kappa(\mathrm{Vir}_2)
\;=\;
12 + 1
\;=\;
13,
\end{equation}
the universal Virasoro complementarity sum
$\kappa(\mathrm{Vir}_c) + \kappa(\mathrm{Vir}_{26-c})
= c/2 + (26-c)/2 = 13$, not~$0$ (Koszul conductor $K = 13$ for
the class-$\mathbf{M}$ Virasoro lane, distinct from the
class-$\mathbf{G}$/$\mathbf{L}$ conductor $K = 0$). This equation
identifies the Virasoro-sector Verdier partner; it does not compute
the full algebra $(V^\natural)^!$.

%% ================================================================
%% Section 4. Moonshine and the shadow obstruction tower
%% ================================================================
\section{Moonshine and the shadow obstruction tower}

\begin{remark}[The $j$-function and the shadow tower]%
\label{rem:moonshine-j-function}%
\index{Monster module!$j$-function and shadow tower}%
\index{$j$-function!shadow tower connection}%
The partition function of $V^\natural$ is the Hauptmodul
$J(\tau) = j(\tau) - 744 = q^{-1} + 196884\,q
+ 21493760\,q^2 + \cdots$,
whose coefficients decompose into Monster representations
(Conway--Norton monstrous moonshine, proved by
Borcherds \cite{Bor92}).

The shadow obstruction tower at the Virasoro scalar level sees
$\kappa = 12$ and the Virasoro shadow coefficients
$S_r$ at $c = 24$; it does not directly encode the
$j$-function coefficients $196884$, $21493760$, \ldots.
The moonshine data live first in the graded Monster module
$V^\natural$ itself: the graded traces of $g \in \mathbb{M}$
recover the McKay--Thompson series $T_g(\tau)$. A corresponding
equivariant shadow-tower statement would require computing the
Monster action on $H^*(\barB(V^\natural))$ and comparing its graded
traces with these series; this is an additional equivariant
Euler-characteristic problem for the bar complex.
The Rademacher exact formula (the genus-$0$ Hauptmodul
property: $J$ is determined by its polar part $q^{-1}$)
is a modular bootstrap constraint on the genus-$0$ partition function,
not a consequence of the degree-$2$ scalar shadow alone.
\end{remark}

\begin{remark}[Monster, BKM, K3, and Yangian lanes]%
\label{rem:moonshine-taxonomy-firewall}%
\index{Monster module!taxonomy firewall}%
\index{Borcherds--Kac--Moody algebra!separated from Monster VOA}%
\index{Yangian!separated from moonshine shadow lane}%
The following objects share coefficients and names but not ambient
categories.
\begin{center}
\begin{tabular}{p{3.2cm}p{8.9cm}}
\textbf{Monster VOA} &
 $V^\natural$ is a holomorphic vertex algebra with
 $\Aut(V^\natural)=\mathbb{M}$; this chapter uses its Virasoro
 scalar projection for $\kappa$, $\Delta_T$, and class~M.\\[2mm]
\textbf{Monster BKM} &
 Borcherds' monster Lie algebra is a generalized
 Kac--Moody algebra whose denominator exponents are coefficients of
 $J$. It is not $B(V^\natural)$, $(V^\natural)^i$, or
 $(V^\natural)^!$.\\[2mm]
\textbf{Igusa/$\Delta_5$ K3 lane} &
 The Igusa/Borcherds product attached to K3 or $\Delta_5$ belongs to
 the Siegel automorphic and BKM denominator lane; its weight and
 Fourier--Jacobi exponents are not the Monster Virasoro scalar
 invariants.\\[2mm]
\textbf{K3 Hall--Drinfeld lane} &
 K3 Hall or Hall--Drinfeld constructions use CoHA/Hall algebra
 multiplication and Mukai data. Their Hochschild or centre statements
 require that Hall ambient.\\[2mm]
\textbf{Drinfeld Yangian lane} &
 Drinfeld Yangians are ordered $E_1$ RTT/Hall--Drinfeld quantum-group
 objects. They have no Virasoro central charge, no Virasoro
 $S_4$, and no scalar conductor value in the
 $\mathbf{G}/\mathbf{L}/\mathbf{C}/\mathbf{M}$ table.
\end{tabular}
\end{center}
Thus a moonshine, BKM, K3, Hall, or Yangian row becomes a shadow
statement only after a named projection to a Virasoro, specified
$\mathcal W$, lattice, or ordered RTT line has been supplied.
\end{remark}

\begin{remark}[The orbifold and shadow class transition]%
\label{rem:moonshine-orbifold-class-transition}%
\index{Monster module!orbifold class transition}%
The $\bZ/2$ orbifold
$V_{\Lambda_{24}} \longrightarrow V^\natural$
simultaneously transforms:
\begin{enumerate}[label=\textup{(\roman*)}]
\item $\dim V_1$: from $24$ to~$0$
 (all weight-$1$ currents killed);
\item $\kappa$: from $24$ to $12$
 (rank-$24$ current contribution absent);
\item shadow class on the scalar projection: from $\mathbf{G}$ to
 $\mathbf{M}$ (finite shadow depth to infinite shadow depth);
\item $\Delta$: from $0$ to $20/71$
 (irreducible shadow metric);
\item $V_2$ structure: $\dim V_2 = 196884$ for both
 $V_{\Lambda_{24}}$ and~$V^\natural$, but the
 decomposition changes from $324$ Heisenberg descendants
 $+\;196560$ lattice states to $1$ Virasoro descendant
 $+\;196883$ Griess primaries carrying the Monster action.
\end{enumerate}
Within the Niemeier--Monster atlas recorded here, this finite
orbifold changes the shadow class. The mechanism:
killing the Heisenberg generators removes the lattice
contribution to $\kappa$, halving it from $24$ to $12$;
the surviving Virasoro sector with
$S_4 = 5/1704 \neq 0$ produces
$\Delta = 20/71 \neq 0$, which the single-line dichotomy
converts to infinite shadow depth.
\end{remark}

\section{Schellekens 71 and the shadow partition}
\label{sec:schellekens-three-mechanism}

\begin{theorem}[Schellekens structured decomposition at
\texorpdfstring{$c=24$}{c=24}]
\label{thm:schellekens-three-mechanism}
\ClaimStatusProvedHere
Let $\mathcal{S}_{24}$ be the $71$ strongly rational holomorphic
central-charge-$24$ rows in Schellekens's list. As a set of rows,
\[
\mathcal{S}_{24}
= \mathcal{S}_{24}^{\mathrm{Niem}}
\sqcup \{V^\natural\}
\sqcup \mathcal{S}_{24}^{\mathrm{nonlat}},
\qquad
71 = 24_{\mathrm{Niemeier}} + 1_{V_1=0} + 46_{\mathrm{nonlat}}.
\]
The Niemeier summand has scalar lattice data
$(\kappa,\mathrm{class},\Delta)=(24,\mathbf{G},0)$. The singleton
$V_1=0$ row, represented by $V^\natural$, has Virasoro scalar data
$(\kappa,\mathrm{class},\Delta)=(12,\mathbf{M},20/71)$.
The remaining $46$ rows are non-lattice Schellekens slots; the row
partition supplies no BV differential, no fixed-sublattice cocycle,
and no Dijkgraaf--Witten vanishing datum for them.
\end{theorem}

\begin{proof}
The row decomposition is Theorem~\ref{thm:schellekens-structured-subset}:
the Niemeier rows are the abelian row together with the twenty-three
level-$1$ simply-laced ADE rows of total rank~$24$, the singleton row is
characterized by $V_1=0$, and the remaining displayed Schellekens rows
form the non-lattice complement. The lattice scalar data are
Theorem~\ref{thm:lattice:niemeier-shadow-universality}. The moonshine
scalar data are Proposition~\ref{prop:moonshine-kappa} and
Equation~\eqref{eq:moonshine-discriminant}. These three inputs give the
displayed partition and its two scalar shadow profiles.
\end{proof}

\begin{remark}[Leech-orbifold presentations and BV data]
M\"oller--Scheithauer generalized-deep-hole presentations are
classification data for VOAs. A BV lift is stronger data: one must
specify, row by row, the cyclic group, the equivariant cocycle on the
relevant central extension of the lattice, the fixed sublattice, the
twisted Jacobi identity, and a BV complex in which
$(Q+[\pi_{\mathrm{orb}},-])^2=0$. The structured row partition therefore
does not imply a uniform chain-level $3$d HT-QFT mechanism.
\end{remark}
